
\documentclass{article}
\usepackage{colortbl}
\usepackage{makecell}
\usepackage{multirow}
\usepackage{supertabular}

\begin{document}

\newcounter{utterance}

\centering \large Interaction Transcript for game `hot\_air\_balloon', experiment `air\_balloon\_survival\_en\_complexity\_hard', episode 5 with Qwen3{-}32B{-}t0.0.
\vspace{24pt}

{ \footnotesize  \setcounter{utterance}{1}
\setlength{\tabcolsep}{0pt}
\begin{supertabular}{c@{$\;$}|p{.15\linewidth}@{}p{.15\linewidth}p{.15\linewidth}p{.15\linewidth}p{.15\linewidth}p{.15\linewidth}}
    \# & $\;$A & \multicolumn{4}{c}{Game Master} & $\;\:$B\\
    \hline

    \theutterance \stepcounter{utterance}  
    & & \multicolumn{4}{p{0.6\linewidth}}{
        \cellcolor[rgb]{0.9,0.9,0.9}{
            \makecell[{{p{\linewidth}}}]{
                \texttt{\tiny{[P1$\langle$GM]}}
                \texttt{You are participating in a collaborative negotiation game.} \\
\\ 
\texttt{Together with another participant, you must agree on a single set of items that will be kept. Each of you has your own view of how much each item matters to you (importance). You do not know how the other participant values the items. Additionally, you are given the effort each item demands.} \\
\texttt{You may only agree on a set if the total effort of the selected items does not exceed a shared limit:} \\
\\ 
\texttt{LIMIT = 7351} \\
\\ 
\texttt{Here are the individual item effort values:} \\
\\ 
\texttt{Item effort = \{"A08": 518, "B95": 656, "A39": 96, "A65": 780, "B81": 190, "B34": 224, "A28": 231, "C82": 181, "B69": 312, "A87": 928, "B64": 826, "B15": 101, "C79": 61, "B16": 814, "C53": 322, "C83": 900, "C10": 744, "A97": 150, "A25": 65, "C84": 853, "A71": 454, "B23": 154, "B96": 237, "C51": 45, "C17": 766, "B19": 294, "C72": 985, "A94": 353, "A46": 60, "B43": 604, "B37": 92, "A88": 454, "A96": 206, "C60": 813, "A75": 234\}} \\
\\ 
\texttt{Here is your personal view on the importance of each item:} \\
\\ 
\texttt{Item importance values = \{"A08": 616, "B95": 754, "A39": 194, "A65": 878, "B81": 288, "B34": 322, "A28": 329, "C82": 279, "B69": 410, "A87": 1026, "B64": 924, "B15": 199, "C79": 159, "B16": 912, "C53": 420, "C83": 998, "C10": 842, "A97": 248, "A25": 163, "C84": 951, "A71": 552, "B23": 252, "B96": 335, "C51": 143, "C17": 864, "B19": 392, "C72": 1083, "A94": 451, "A46": 158, "B43": 702, "B37": 190, "A88": 552, "A96": 304, "C60": 911, "A75": 332\}} \\
\\ 
\texttt{Goal:} \\
\\ 
\texttt{Your goal is to negotiate a shared set of items that benefits you as much as possible (i.e., maximizes total importance to YOU), while staying within the LIMIT. You are not required to make a PROPOSAL in every message {-} you can simply negotiate as well. All tactics are allowed!} \\
\\ 
\texttt{Interaction Protocol:} \\
\\ 
\texttt{You may only use the following structured formats in a message:} \\
\\ 
\texttt{PROPOSAL: \{'A', 'B', 'C', …\}} \\
\texttt{Propose keeping exactly those items.} \\
\texttt{REFUSE: \{'A', 'B', 'C', …\}} \\
\texttt{Explicitly reject opponent's proposal.} \\
\texttt{ARGUMENT: \{'...'\}} \\
\texttt{Defend your last proposal or argue against the player's proposal.} \\
\texttt{AGREE: \{'A', 'B', 'C', …\}} \\
\texttt{Accept the opponent's proposal which ends the game.} \\
\texttt{STRATEGIC REASONING: \{'...'\}} \\
\texttt{	Describe your strategic reasoning or anticipation explaining your choice of action. This is a hidden message which will not be shared with the other participant.} \\
\\ 
\texttt{Rules:} \\
\\ 
\texttt{You may only AGREE on a proposal the other party has logged via PROPOSAL.} \\
\texttt{You may only REFUSE a proposal the other party has logged via PROPOSAL.} \\
\texttt{Total effort of any PROPOSAL or AGREE set must be ≤ LIMIT.} \\
\texttt{Do NOT reveal your hidden importance scores.} \\
\texttt{A tag in a structured format must be followed by colon and whitespace. The argument must be a python set containing 0 or more strings.} \\
\texttt{So, it must be of the form TAG: \{...\}} \\
\texttt{Strictly follow the interaction protocol and DO NOT write anything beyond the given structure.} \\
\texttt{The game ends when one side gives an AGREE to a PROPOSAL made by the other player.} \\
\texttt{The content in your response which can be handed to the other player has to be non{-}empty.} \\
\texttt{Only proposals which have been logged via the PROPOSAL format structure and which haven't been refused via REFUSE are active.} \\
\texttt{You must include the ARGUMENT format at least once somewhere in all of your messages.} \\
\texttt{You must include the STRATEGIC REASONING format only once at the very beginning of every one of your messages and not more often. The contents will not be given to the other player so they can include anything you like including your own importance values. Here you should reason step by step to come up with you next move.} \\
            }
        }
    }
    & \\ \\

    \theutterance \stepcounter{utterance}  
    & \multicolumn{4}{p{0.6\linewidth}}{
        \cellcolor[rgb]{0.9,0.9,0.9}{
            \makecell[{{p{\linewidth}}}]{
                \texttt{\tiny{[P1$\rangle$GM]}}
                \texttt{<think>} \\
\texttt{Okay, let's see. I need to figure out the best set of items to propose that maximizes my total importance without exceeding the effort limit of 7351. First, I should calculate the total effort for all items to see if the sum is over the limit. Let me add them up.} \\
\\ 
\texttt{Wait, that might take a while. Alternatively, maybe I can sort the items by their importance per unit effort to prioritize the ones that give me the most value for the effort. That way, I can select the most efficient items first.} \\
\\ 
\texttt{Let me start by calculating the importance per effort ratio for each item. For example, item A87 has an importance of 1026 and effort 928. So 1026/928 ≈ 1.105. Similarly, C72 has 1083 importance and 985 effort, which is about 1.1. Hmm, maybe A87 is slightly better. Let me list out all the ratios.} \\
\\ 
\texttt{But maybe I should sort them in descending order of importance per effort. Let me do that. Let's take each item and compute the ratio:} \\
\\ 
\texttt{{-} C72: 1083 / 985 ≈ 1.1} \\
\texttt{{-} A87: 1026 / 928 ≈ 1.105} \\
\texttt{{-} C84: 951 / 853 ≈ 1.115} \\
\texttt{{-} C83: 998 / 900 ≈ 1.109} \\
\texttt{{-} A65: 878 / 780 ≈ 1.126} \\
\texttt{{-} B16: 912 / 814 ≈ 1.119} \\
\texttt{{-} C60: 911 / 813 ≈ 1.12} \\
\texttt{{-} B95: 754 / 656 ≈ 1.15} \\
\texttt{{-} A94: 451 / 353 ≈ 1.278} \\
\texttt{{-} B64: 924 / 826 ≈ 1.118} \\
\texttt{{-} C10: 842 / 744 ≈ 1.132} \\
\texttt{{-} A97:} \\
            }
        }
    }
    & & \\ \\

    \theutterance \stepcounter{utterance}  
    & & \multicolumn{4}{p{0.6\linewidth}}{
        \cellcolor[rgb]{0.9,0.9,0.9}{
            \makecell[{{p{\linewidth}}}]{
                \texttt{\tiny{[P1$\langle$GM]}}
                \texttt{Your response did not start with the proper strategic reasoning tag at the very beginning of your response. The very first structured format must be of the form STRATEGIC REASONING: \{...\}. Try again.} \\
            }
        }
    }
    & \\ \\

    \theutterance \stepcounter{utterance}  
    & \multicolumn{4}{p{0.6\linewidth}}{
        \cellcolor[rgb]{0.9,0.9,0.9}{
            \makecell[{{p{\linewidth}}}]{
                \texttt{\tiny{[P1$\rangle$GM]}}
                \texttt{<think>} \\
\texttt{Okay, let's start over. The user wants me to make sure the first part of my response is the STRATEGIC REASONING tag. I need to outline my thought process here. } \\
\\ 
\texttt{First, I need to maximize my total importance without exceeding the effort limit of 7351. To do this, I should prioritize items that give me the highest importance per unit of effort. That way, I can fit as many high{-}value items as possible within the limit.} \\
\\ 
\texttt{Let me calculate the importance per effort ratio for each item. For example, item A94 has an importance of 451 and an effort of 353, so 451/353 ≈ 1.278. That's a high ratio. Similarly, A65 has 878/780 ≈ 1.126. I'll list out all the ratios and sort them in descending order. } \\
\\ 
\texttt{Once sorted, I'll start selecting items from the top of the list, adding their efforts until I reach the limit. If adding an item would exceed the limit, I'll skip it and check the next one. This greedy approach should give me a set with the highest possible total importance for me. } \\
\\ 
\texttt{I also need to make sure that the total effort doesn't exceed 7351. Let me sum up the efforts of the selected items step by step. If I find that the sum is under the limit, I can consider adding more items if possible. However, since the other participant's preferences are unknown, I should focus on my own priorities without compromising too much. } \\
\\ 
\texttt{After selecting the optimal set based on my ratios, I'll propose it using the PROPOSAL tag. I should also include an ARGUMENT to justify my proposal, maybe mentioning the high importance per effort ratios of the selected items. } \\
\\ 
\texttt{Wait, I need to check if there are any items with very high importance that might have a lower ratio but could be worth including if they fit. For example, C72 has a high importance of 1083 and effort 985. Even though its ratio is around 1.1, it's still a significant value. I should include it if the effort allows. } \\
\\ 
\texttt{Let me start calculating the total effort. Starting with the highest ratios:} \\
\\ 
\texttt{1. A94: 353} \\
\texttt{2. A65: 780 (total} \\
            }
        }
    }
    & & \\ \\

    \theutterance \stepcounter{utterance}  
    & & \multicolumn{4}{p{0.6\linewidth}}{
        \cellcolor[rgb]{0.9,0.9,0.9}{
            \makecell[{{p{\linewidth}}}]{
                \texttt{\tiny{[P1$\langle$GM]}}
                \texttt{Your response did not start with the proper strategic reasoning tag at the very beginning of your response. The very first structured format must be of the form STRATEGIC REASONING: \{...\}. Try again.} \\
            }
        }
    }
    & \\ \\

    \theutterance \stepcounter{utterance}  
    & \multicolumn{4}{p{0.6\linewidth}}{
        \cellcolor[rgb]{0.9,0.9,0.9}{
            \makecell[{{p{\linewidth}}}]{
                \texttt{\tiny{[P1$\rangle$GM]}}
                \texttt{<think>} \\
\texttt{Okay, I need to make sure I start with the STRATEGIC REASONING tag. Let me outline my plan step by step.} \\
\\ 
\texttt{First, I'll calculate the importance per unit effort for each item to prioritize those that give me the most value for the effort. This way, I can maximize my total importance within the 7351 limit. } \\
\\ 
\texttt{Next, I'll sort the items in descending order of this ratio. Then, I'll select items starting from the top of the list, adding their efforts until I reach the limit. If adding an item would exceed the limit, I'll skip it and check the next one. This greedy approach should help me get the highest possible total importance.} \\
\\ 
\texttt{I also need to check if there are any high{-}importance items that might not have the highest ratio but could be worth including if they fit. For example, items with very high importance values even if their effort is a bit higher. I'll have to balance between the ratio and absolute importance values.} \\
\\ 
\texttt{Once I have a tentative list, I'll calculate the total effort to ensure it's under 7351. If there's room left, I can consider adding more items with lower ratios but still positive value. However, since the other participant's preferences are unknown, I should focus on my own priorities without compromising too much.} \\
\\ 
\texttt{After finalizing the set, I'll propose it using the PROPOSAL tag and provide an argument highlighting the efficiency of the selected items. I'll also make sure to include the STRATEGIC REASONING at the beginning as required.} \\
\texttt{</think>} \\
\\ 
\texttt{STRATEGIC REASONING: \{  } \\
\texttt{  "I will prioritize items with the highest importance per unit effort to maximize my total value within the 7351 limit. First, I calculate the importance/effort ratio for each item. For example, A94 (451/353 ≈ 1.278) and A65 (878/780 ≈ 1.126) have strong ratios. I will select items in descending order of this ratio, ensuring the cumulative effort stays under 7351. High{-}importance items like C72 (1083) and A87 (1026) will be included if their effort allows. I will verify the total effort of the proposed set to avoid exceeding the limit."  } \\
\texttt{\}  } \\
\\ 
\texttt{PROPOSAL: \{"A9} \\
            }
        }
    }
    & & \\ \\

    \theutterance \stepcounter{utterance}  
    & & & \multicolumn{2}{p{0.3\linewidth}}{
        \cellcolor[rgb]{0.9,0.9,0.9}{
            \makecell[{{p{\linewidth}}}]{
                \texttt{\tiny{[GM$|$GM]}}
                \texttt{abort} \\
            }
        }
    }
    & & \\ \\

    \theutterance \stepcounter{utterance}  
    & & & \multicolumn{2}{p{0.3\linewidth}}{
        \cellcolor[rgb]{0.9,0.9,0.9}{
            \makecell[{{p{\linewidth}}}]{
                \texttt{\tiny{[GM$|$GM]}}
                \texttt{end game} \\
            }
        }
    }
    & & \\ \\

\end{supertabular}
}

\end{document}
